\documentclass[11pt,a4paper]{article}

\usepackage[utf8]{inputenc}
\usepackage[spanish]{babel}
\usepackage{amsmath, amsfonts, amssymb}
\usepackage{geometry}
\usepackage{enumitem}

\geometry{top=2.5cm, bottom=2.5cm, left=2.5cm, right=2.5cm}

\begin{document}

\begin{center}
    \textbf{\Large Probabilidad y Estadística} \\
    \vspace{0.2cm}
    \textit{Universidad Tecnológica Nacional}
\end{center}

\vspace{0.5cm}

\section*{Ejercicio 76: Probabilidad Condicional y Reducción del Espacio}

\textbf{Enunciado:} De 100 personas que solicitaron empleo en una firma, 40 tenían experiencia anterior, y 30 tenían un certificado profesional[cite: 6]. Sin embargo, 20 de los solicitantes contaban con ambos antecedentes[cite: 6]. Determine la probabilidad condicional de que un solicitante aleatoriamente elegido tenga un certificado, dado que tiene alguna experiencia anterior[cite: 6].

\vspace{0.5cm}
\hrule
\vspace{0.5cm}

\subsection*{Resolución paso a paso}

Este problema nos permite observar cómo la información adicional (una condición) altera el tamaño de nuestro espacio muestral. Lo abordaremos desde dos enfoques equivalentes para consolidar el concepto.

\subsubsection*{1. Definición de Sucesos y Extracción de Datos}
Primero, definimos los sucesos de interés y anotamos las cardinalidades (cantidad de casos) provistas por el enunciado:
\begin{itemize}
    \item $\Omega$: Conjunto total de solicitantes de empleo ($|\Omega| = 100$).
    \item $E$: El solicitante tiene experiencia anterior ($|E| = 40$).
    \item $C$: El solicitante tiene un certificado profesional ($|C| = 30$).
    \item $E \cap C$: El solicitante posee ambos antecedentes ($|E \cap C| = 20$).
\end{itemize}

El objetivo es calcular $P(C \mid E)$, es decir, la probabilidad de que el candidato tenga certificado \textbf{dado} que ya sabemos como un hecho cierto que tiene experiencia.

\subsubsection*{2. Enfoque Intuitivo: Reducción del Espacio Muestral}
Al saber que el evento $E$ ha ocurrido, nuestro universo de análisis descarta a las 60 personas sin experiencia. El nuevo espacio muestral se reduce estrictamente al conjunto $E$.

\begin{itemize}
    \item \textbf{Casos Posibles (Nuevo Universo):} Son las personas que cumplen la condición de tener experiencia ($|E| = 40$).
    \item \textbf{Casos Favorables:} Dentro de ese nuevo universo de 40 personas, las únicas que tienen certificado son las que pertenecen a la intersección de ambos conjuntos ($|E \cap C| = 20$).
\end{itemize}

Aplicando la regla de Laplace sobre este universo restringido:
$$ P(C \mid E) = \frac{\text{Casos Favorables dentro de } E}{\text{Casos Totales en } E} = \frac{|E \cap C|}{|E|} = \frac{20}{40} = 0,5 $$

\subsubsection*{3. Enfoque Formal: Definición Axiomática}
El mismo resultado se obtiene aplicando la fórmula general de probabilidad condicional sobre las probabilidades marginales del espacio muestral original ($\Omega$).

Calculamos las probabilidades de los sucesos dividiendo por el total histórico de 100 solicitantes:
$$ P(E) = \frac{40}{100} = 0,40 $$
$$ P(E \cap C) = \frac{20}{100} = 0,20 $$

Finalmente, reemplazamos en la definición de la probabilidad condicional:
$$ P(C \mid E) = \frac{P(E \cap C)}{P(E)} = \frac{0,20}{0,40} = 0,5 $$

\vspace{0.3cm}
\noindent \textbf{Conclusión:} Al poseer el dato de que el candidato tiene experiencia previa, la probabilidad de que cuente con un certificado profesional es del $50\%$.

\end{document}
